% !TeX program = lualatex
% Copyright (C) 2026 IEEE UFG Computer Society contributors.
% Distributed under the LaTeX Project Public License 1.3c or later.
% See LICENSE and development/MANIFEST.md.
%
% MODELO DE ATA DE REUNIÃO
% Substitua todos os campos entre colchetes. Registre fatos, decisões e
% responsáveis; remova itens opcionais que não se aplicarem à reunião.
% Para reuniões com muitos participantes, anexe uma lista de presença.
\documentclass[
  paper-size=a4,
  font-size=11pt,
  two-sided=false,
  title-page=false,
  section-numbering=plain,
  inner-margin=22mm,
  outer-margin=22mm,
  top-margin=24mm,
  bottom-margin=28mm,
  color-model=rgb
]{ieeeufgcs}
\usepackage[brazilian]{babel}
\usepackage{indentfirst}
\raggedbottom

% Edite estes valores uma vez; eles aparecem no título e na identificação.
\newcommand{\AtaNumero}{[número/ano]}
\newcommand{\AtaData}{[dia de mês de ano]}

\IEEECSSetup{
  document-type={Ata de reunião},
  title={Ata nº \AtaNumero},
  subtitle={Diretoria Executiva},
  running-title={Ata nº \AtaNumero},
  institution={Capítulo Estudantil da IEEE Computer Society na UFG},
  author={Secretaria},
  date={\AtaData},
  location={Goiânia},
  identifier={Minuta para aprovação}
}

\begin{document}
\maketitle

\begin{IEEECSBrandedTable}{
  >{\raggedright\arraybackslash}m{0.28\linewidth}
  >{\raggedright\arraybackslash}X
}
\IEEECSHeaderCell{Identificação} & \IEEECSHeaderCell{Registro} \\
Órgão ou grupo &
Capítulo Estudantil da IEEE Computer Society na UFG --- Diretoria Executiva.
\\
Tipo de reunião & [Ordinária, extraordinária ou outra]. \\
Data & \AtaData. \\
Horário & Início: [hh:mm]. Encerramento: [hh:mm]. \\
Local ou modalidade & [Sala e unidade, endereço ou plataforma virtual]. \\
Presidência da reunião & [Nome e função]. \\
Secretaria da reunião & [Nome e função]. \\
Quórum, se aplicável & [Regra adotada e presença verificada]. \\
\end{IEEECSBrandedTable}

\section*{Participantes}

\textbf{Presentes:} [Nomes completos e funções, quando relevantes].

\textbf{Ausências justificadas:} [Nomes, se houver].

\textbf{Convidados:} [Nomes e vínculos, se houver].

% Remova a linha de quórum se não houver regra aplicável.
% Anexe lista de presença se a relação for extensa ou exigir assinatura.

\section*{Abertura e pauta}

A reunião foi aberta por [nome] às [hh:mm]. A pauta comunicada aos
participantes compreendeu:
\begin{enumerate}
  \item \emph{[Primeiro assunto]};
  \item \emph{[Segundo assunto]}.
\end{enumerate}

\section*{Discussões e deliberações}

\subsection*{Item 1 --- [assunto]}

[Registre os fatos e os argumentos necessários para compreender a decisão,
sem transcrever falas nem incluir dados pessoais desnecessários.]

\textbf{Deliberação:} [Descreva a decisão ou registre que não houve decisão].
\textbf{Votação, se aplicável:} [Resultado, abstenções e impedimentos].

\subsection*{Item 2 --- [assunto]}

[Registre a síntese da discussão e o resultado do item. Duplique esta
subseção para os demais assuntos efetivamente tratados.]

\textbf{Deliberação:} [Decisão, adiamento ou ausência de decisão].

\section*{Encaminhamentos}

\begin{IEEECSBrandedTable}{
  >{\raggedright\arraybackslash}X
  >{\raggedright\arraybackslash}m{0.24\linewidth}
  >{\raggedright\arraybackslash}m{0.17\linewidth}
}
\IEEECSHeaderCell{Ação ou entrega} &
\IEEECSHeaderCell{Responsável} &
\IEEECSHeaderCell{Prazo} \\
\emph{[Ação aprovada e resultado esperado]} &
\emph{[Nome ou função]} & \emph{[Data]} \\
\emph{[Outra ação, se houver]} &
\emph{[Nome ou função]} & \emph{[Data]} \\
\end{IEEECSBrandedTable}

\section*{Encerramento e aprovação}

Nada mais havendo a tratar, [nome] encerrou a reunião às [hh:mm].
Esta ata foi redigida por [nome] e será [forma de leitura, revisão e
aprovação adotada pelo grupo]. [Registre a data da aprovação após ela
ocorrer.]

\begin{IEEECSBrandedTable}{
  >{\raggedright\arraybackslash}m{0.42\linewidth}
  >{\raggedright\arraybackslash}X
}
\IEEECSHeaderCell{Nome e função} &
\IEEECSHeaderCell{Assinatura e data} \\
\emph{[Nome de quem presidiu]} --- Presidência &
\emph{[Assinatura e data]} \\
\emph{[Nome de quem secretariou]} --- Secretaria &
\emph{[Assinatura e data]} \\
\end{IEEECSBrandedTable}

% Anexos opcionais: lista de presença, material apresentado, votação ou
% documento aprovado. Identifique-os aqui e guarde-os com a ata final.

\end{document}
